% ####################################################################
% Basic configuration and packages
% ####################################################################
% Package for discovering wrong and outdated usage of LaTeX.
% More information to be found in l2tabu English version.
\RequirePackage[l2tabu, orthodox]{nag}
% Class of LaTeX document: {size of paper, size of font}[document class]
\documentclass[a4paper,11pt]{article}



% ######################################
% Packages not tied to particular normal language
% ######################################
% This package should improved spaces in the text
\usepackage{microtype}
% Add few important symbols, like text Celcius degree
\usepackage{textcomp}



% ######################################
% Polonization of LaTeX document
% ######################################
% Basic polonization of the text
\usepackage[MeX]{polski}
% Switching on UTF-8 encoding
\usepackage[utf8]{inputenc}
% Adding font Latin Modern
\usepackage{lmodern}
% Package is need for fonts Latin Modern
\usepackage[T1]{fontenc}



% ------------------------------------------------------
% Setting margins
% ------------------------------------------------------
\usepackage[a4paper, total={14cm, 25cm}]{geometry}



% ------------------------------------------------------
% Setting vertical spaces in the text
% ------------------------------------------------------
% Setting space between lines
\renewcommand{\baselinestretch}{1.1}

% Setting space between lines in tables
\renewcommand{\arraystretch}{1.4}



% ------------------------------------------------------
% Packages for scientific papers
% ------------------------------------------------------
% Switching off \lll symbol, that I guess is representing letter "Ł"
% It collide with `amsmath' package's command with the same name
\let\lll\undefined
% Basic package from American Mathematical Society (AMS)
\usepackage[intlimits]{amsmath}
% Equations are numbered separately in every section
\numberwithin{equation}{section}

% Other very useful packages from AMS
\usepackage{amsfonts}
\usepackage{amssymb}
\usepackage{amscd}
\usepackage{amsthm}

% Package with better looking calligraphy fonts
\usepackage{calrsfs}

% Package with better looking greek letters
% Example of use: pi -> \uppi
\usepackage{upgreek}
% Improving look of lambda letter
\let\oldlambda\Lambda
\renewcommand{\lambda}{\uplambda}




% ------------------------------------------------------
% BibLaTeX
% ------------------------------------------------------
% Package biblatex, with biber as its backend, allow us to handle
% bibliography entries that use Unicode symbols outside ASCII
\usepackage[
language=polish,
backend=biber,
style=alphabetic,
url=false,
eprint=true,
]{biblatex}

\addbibresource{Rachunek-prawdopodobieństwa-i-statystyka-Bibliography.bib}





% ------------------------------------------------------
% Defining new environments (?)
% ------------------------------------------------------
% Defining enviroment "Wniosek"
\newtheorem{corollary}{Wniosek}
\newtheorem{definition}{Definicja}
\newtheorem{theorem}{Twierdzenie}





% ------------------------------------------------------
% Local packages
% You need to put them in the same directory as .tex file
% ------------------------------------------------------
% Package containing various command useful for working with a text
\usepackage{./Local-packages/general-commands}
% Package containing commands and other code useful for working with
% mathematical text
\usepackage{./Local-packages/math-commands}



% ######################################
% Local macros
% ######################################
\newcommand{\Orzel}{\text{O}}
\newcommand{\Reszka}{\text{R}}





% ######################################
% Package "hyperref"
% They advised to put it on the end of preambule
% ######################################
% It allows you to use hyperlinks in the text
\usepackage{hyperref}










% ####################################################################
% Title and author of the text
\title{Rachunek prawdopodobieństwa \\
  {\Large Błędy i~uwagi}}

\author{Kamil Ziemian}


% \date{}
% ####################################################################










% ####################################################################
\begin{document}
% ####################################################################





% ######################################
\maketitle % Tytuł całego tekstu
% ######################################





% ######################################
\section{Patrick Billingsley
  \textit{Prawdopodobieństwo i~miara},
  \cite{BillingsleyPrawdopodobienstwoIMiara2009}}
% ############################

\vspace{0em}


% ##################
\CenterBoldFont{Uwagi}

\vspace{0em}


\noindent
Warto byłoby przepisać tę~książkę w~\LaTeX u. Widać bowiem, że~wzory
matematyczne zostały zapisane w~jakimś, zapewne pochodzącym z~czasów PRLu
systemie i~symbole są przez to~niskiej jakości, takie pikselowate.

\VerSpaceFour





\noindent
Pionowe kreski w~wykresach funkcji schodkowych, takich jak
te~na~stronie 17, powinny być rysowane linią przerywaną, nie zaś jak
obecnie ciągłą. Stosuje~się to oczywiście do wszystkich funkcji, które
posiadają nieciągłość w~formie skoku.

\VerSpaceFour





% ##################
\CenterBoldFont{Uwagi do~konkretnych stron}


\noindent
\Str{17} Przez liczbę dwójkowo-wymierną rozumie chyba Billingsley liczbę
którą da~się zapisać w~postaci $p / 2^{ n }$, $p < 2^{ n }$.

\VerSpaceFour





\noindent
\Str{17} Pojęcie rzędu przedziału może stwarzać problem, bo~wydaje~się,
że~Billingsley przyjął, iż~liczby $a_{ 0 }, a_{ 1 }, \ldots, a_{ n }$ tworzą
podział rzędu $n$, jeśli każda z~nich da~się zapisać w~postaci
$p / 2^{ n }$, niezależnie od~tego, czy jest to forma nieskracalna czy nie.
Wnioskuję to z~tego, że~ciąg $0, 1/4, 2/4, 3/4, 1$ powinien tworzyć podział
rzędu 2, jednak $2/4 = 1/2$.

Należałoby więc rząd przedziału zdefiniować w~ten sposób, że~każdą
z~liczb $a_{ i }$ można zapisać w~postaci $p_{ i } / 2^{ n }$ przy
czym co~najmniej inna z~nich jest w~postaci nieskracalnej. Tym samym
pojęcie to~stanie~się jednoznaczne.

\VerSpaceFour





% ##################
% \CenterBoldFont{Błędy}


% \begin{center}

%   \begin{tabular}{|c|c|c|c|c|}
%     \hline
%     Strona & \multicolumn{2}{c|}{Wiersz} & Jest & Powinno być \\
%     \cline{2-3}
%     & Od góry & Od dołu & & \\
%     \hline
    %     & & & & \\
    %     & & & & \\
    %     & & & & \\
    %     & & & & \\
    %     & & & & \\
    %     & & & & \\
    %     & & & & \\
    %     & & & & \\
    %     & & & & \\
    %     & & & & \\
    %     & & & & \\
    %     & & & & \\
    %     & & & & \\
    %     & & & & \\
    %     & & & & \\
    %     & & & & \\
%     \hline
%   \end{tabular}

% \end{center}

% \VerSpaceSix


% ############################









% ####################################################################
\section{William Feller \textit{Wstęp do rachunku prawdopodobieństwa. Tom~I},
  \cite{FellerWstepDoRachunkuPrawdopodobienstwaVolI2006}}

\label{sec:Feller-Wstep-do-rachunku-ETC}
% ####################################################################



% % ######################################
% \subsection{Uwagi~konkretnych stron}

% \label{subsec:Uwagi-do-konkretnych-stron}
% % ######################################




% ######################################
\subsection{Uwagi do~konkretnych stron}

\label{subsec:Uwagi-do-konkretnych-stron}
% ######################################



\noindent
\StrWierszeGora{13}{9--11} Według mnie zdanie „Potrzebujemy ścisłej
matematycznej teorii opartej na zasadach ogólnie przyjętych w~geometrii
i~mechanice.”, powinno być zastąpione przez „Potrzebujemy ścisłej
matematycznej teorii opartej na ogólnych zasadach, jak tych przyjętych
w~geometrii i~mechanice.”.

{\Large Od strony 17 należy zacząć ponowne czytanie książki.}



% ##################
\CenterBoldFont{Błędy}


\begin{center}

  \begin{tabular}{|c|c|c|c|c|}
    \hline
    Strona & \multicolumn{2}{c|}{Wiersz} & Jest
                              & Powinno być \\ \cline{2-3}
    & Od góry & Od dołu & & \\
    \hline
    18 & & 5 & \ldots lub & \ldots{} lub \\
    19 & & & $- |$ & $- \}$ \\
    28 & & 9 & wielu sposobami & na wiele sposobów \\
    % & & & & \\
    % & & & & \\
    % & & & & \\
    \hline
  \end{tabular}

\end{center}

\VerSpaceSix



% ############################










% ##########################################################
\section{Zdzisław Hellwig \textit{Elementy rachunku
    prawdopodobieństwa i~statystyki matematycznej},
  \parencite{Hellwig-Elementy-rachunku-prawdopodobienstwa-ETC-Pub-1987}}

\label{sec:Hellwig-Elementy-Rachunku-ETC}
% ##########################################################



% ################################################
\subsection{Uwagi}

\label{subsec:Uwagi}
% ################################################


W~rachunku prawdopodobieństwa musimy często rozważać ciągi elementów
pewnego zbioru $E$, czyli odwzorowania $a : \{ 1, 2,\ldots, n \} \to E$. By praca
z~nimi była wygodniejsza, przyjmujemy, że~przez $a_{ i }$ będziemy oznaczać
$i$-ty element danego ciągu, zaś sam ciąg będziemy oznaczać przez
$( a_{ i } )$. Natomiast zbiór wartości ciągu będziemy oznaczać, przez
$\{ a_{ i } \}$. W~omawianej książce tego typu ciągi nazywane są
\textit{kombinacjami},
a~dokładniej???? \textit{kombinacjami z~powtórzeniami}.

W~książce spotykamy~się też z~następującą sytuacją. Chcemy mieć ciąg,
powiedzmy, że~dwuelementowy, taki że~$a_{ 1 } \in E_{ 1 }$ i~$a_{ 2 } \in E_{ 2 }$,
przy czym może zachodzić $E_{ 1 } \nsubseteq E_{ 2 }$ i~$E_{ 2 } \nsubseteq E_{ 1 }$.
Sformalizowanie takich ciągów jest możliwe, lecz nie będziemy~się tym
zajmować, gdyż na tym etapie nie przyniesie to wiele korzyści.










% ################################################
\subsection{Uwagi do~konkretnych stron}

\label{subsec:Uwagi-do-konkretnych-stron}
% ################################################



\noindent
\Str{???}


\noindent
\Str{35} Autor rozpoczyna tu dyskusję pojęcia \textbf{zdarzenia losowego}
i~\textbf{zdarzenia elementarnego}, które to są fundamentalne dla~całego
rachunku prawdopodobieństwa. Książka będzie jeszcze wielokrotnie do nich
wracać i~omawiać je pod różnymi kątami, warto jednak zatrzymać~się już na
początku tych rozważań i~możliwie dokładnie omówić ich sens. Nim do tego
przejdziemy, musimy zaznaczyć, że~pojęcie „zdarzenia” i~„zdarzenia losowego”
będziemy traktować jako synonimy, chyba że~powiedziano inaczej.

Nasze rozważania nie będą miały charakteru ściśle formalnego, choć mamy
nadzieję, że~będzie je cechował możliwie wysoki poziom ścisłości. Wybieramy
takie podejście, gdyż naszym celem nie jest po prostu wprowadzenie
matematycznych pojęć, lecz wyjaśnienie czemu je wprowadzamy i~jaki jest ich
ogólny sens. Dlatego też naszym punktem wyjścia jest omówienie problemu,
który przenika cały rachunek prawdopodobieństwa.

Problem ten jest następujący. Mamy doświadczenie, które jest powtarzalne
i~zawsze produkuje pewien konkretnych wynik. Na początek przyjmijmy, że~dane
doświadczenie może zwrócić tylko skończoną liczbę wyników $e_{ 1 }$,
$e_{ 2 }$, \ldots, $e_{ n }$. Jest to uzasadnione tym, że~zrozumienie wariantu
problemu, gdy liczba wyników jest
skończona, samo w~sobie jest ogromnie ważne. Poza tym, uogólnienie tych
rozważań na przypadek, gdy doświadczenie może zwrócić jeden z~nieskończonej
liczby rezultatów jest wyjątkowo trudne.

Był ułatwić nasze rozważania, rozpatrzmy jako model takiego doświadczenie
problem urny, która zawiera $12$ kul. Pięć z~tych kul jest białych i~każda
na swej powierzchni zapisany numer od~$1$ do $5$, pozostałe siedem jest
czarnych i~ponumerowanych liczbami od $1$ do $7$. Poza tym kule można uważać
za identyczne. Doświadczenie będzie polegało na wyciągnięciu jednej kuli
z~urny, bez zaglądania do wnętrza urny,
tak że przed wyciągnięcie kuli, nie znamy ani jej koloru ani numeru.
Doświadczenie to jest powtarzalne, bo zawsze po wyciągnięciu kuli, możemy ją
włożyć z~powrotem do urny. Ten model doświadczenia można by dalej
doprecyzowywać, jednak nie jest to potrzebne, gdyż jest on intuicyjnie prosty
do zrozumienia.

W~omawianym przypadku, zdarzeniem elementarnym jest wyciągnięci jednej kuli
z~urny. Można też powiedzieć, że~to pojęcie zdarzenia elementarnego
powstało po to, by opisać zdarzenia takie jak wyciągnięcie jednej kuli
z~urny. Przy czym słowa „zdarzenie” użyliśmy tu w~normalny sensie, nie
zaś jako synonimu „zdarzenia losowego”. Zauważmy bowiem, że w~problemach
z~rachunku prawdopodobieństwa zawsze lub prawie zawsze mamy do czynienia
z~jakimś doświadczeniem i~dzięki naszemu zrozumieniu
natury tego doświadczenia, jesteśmy w~stanie powiedzieć, co jest zdarzeniem
elementarny. Zaś bez określenia zdarzeń elementarnych nie jesteśmy
w~stanie ani poprawnie sformułować, ani rozwiązań większości problemów
z~rachunku prawdopodobieństwa.

Mówiąc inaczej, w~problemach z~rachunku prawdopodobieństwa prawie zawsze
mamy podane jakieś doświadczenie powtarzalne, takie jak wyciąganie kul
z~urny. Dzięki naszemu zrozumieniu tego doświadczenia, potrafimy powiedzieć,
jaki są możliwe rezultaty jego pojedynczego wykonania. Następnie każdy
możliwy rezultat tego doświadczenia uznajemy za zdarzenie elementarne.
Zauważmy, że~do tej pory nic nie mówiliśmy o~prawdopodobieństwie jaką
posiadają zdarzenia elementarne, na to przyjdzie czas później.

Te rozważania można podsumować mówiąc, że~pojęcie zdarzenia elementarnego
powstało po to, by opisać pojedynczy wynik powtarzalnego doświadczenia.
To utożsamienie, które ciągle spotykamy rozwiązując problemy z~rachunku
prawdopodobieństwa, pozwala nam na podstawie zrozumienie doświadczenia,
zidentyfikować zdarzenia elementarne.

Rozważmy teraz, jakie problemy możemy rozważać w~kontekście powtarzalnego
doświadczenie, takiego jak wyciąganie kul z~naszej urny. By uprościć
rozważania, zdarzenie elementarne
polegające na wyciągnięciu kuli białej o~numerze $1$ oznaczamy $e_{ 1 }$,
kuli o~numerze~$2$ oznaczamy $e_{ 2 }$ i~tak do $e_{ 5 }$. Wyciągnięcie kuli
czarnej o~numerze $1$ oznaczamy przez $e_{ 6 }$, kuli czarnej o~numerze $2$
przez $e_{ 7 }$, etc., aż do $e_{ 12 }$.

Najprostsze problem jaki możemy w~tym kontekście postawić, brzmi następująco.
Wybieramy jedną kulę, niech będzie to kula biała o~numerze $1$,
jej wyciągnięciu odpowiada zdarzenie elementarne $e_{ 1 }$. Wyciągamy kulę
z~urny, czyli zaszło zdarzenie $e_{ j }$, dla $j \in \{ 1, 2, 3, \ldots, 12 \}$. Czy
wyciągnęliśmy białą kulę o~numerze~$1$? Innymi słowy, czy
$e_{ j } = e_{ 1 }$? Należy podkreślić, że~wciąż nic nie powiedzieliśmy
o~prawdopodobieństwie zdarzeń elementarnych, pytamy~się tylko o~to, czy
w~danym eksperymencie zaszło wybrane zdarzenie, czy nie? Mam tu więc do
czynienia tylko z~identyfikacją wyniku doświadczenia, co nie ma nic wspólnego
z~szeroko pojętym prawdopodobieństwem.

Przykład powyższy pokazuje, czemu w~tego typu problemach wygodne jest
oznaczanie tym samym symbolem zarówno kuli, jak i~zdarzenia elementarnego
polegającego na jej wyciągnięciu, co jednak jest bardzo niechlujnym
postępowaniem. Biała kula o~numerze $1$ jest pewną rzeczą, zaś $e_{ 1 }$ jest
zdarzeniem polegającym na wyciągnięciu tej kuli z~urny, a~rzeczy i~zdarzenia
w których rzeczy biorą udział, to zupełnie inne „obiekty”. Niemniej zdarzenie
elementarne $e_{ 1 }$, które polega na wyciągnięciu białej kuli
o~numerze~$1$, jednoznacznie pozwala tę kulę zidentyfikować. Z~tego względu
niektórzy piszą „Rozważmy kulę $e_{ 1 }$\ldots”, co bardziej precyzyjnie
wysłowilibyśmy pisząc „Rozważmy kulę, której wyciągnięcie z~urny jest
zdarzeniem elementarnym oznaczanym przez $e_{ 1 }$\ldots” Pozwolimy sobie
korzystać z~tego typu skrótów myślowych, jeśli uprości to nasze rozważania.

Następny problem jaki możemy postawić w~kontekście tego doświadczenia jest
następujący. Wyciągamy kulę z~urny, czyli zaszło zdarzenie $e_{ j }$, dla
$j \in \{ 1, 2, 3,\ldots, 12 \}$. Czy wyciągnęliśmy białą kulę z~numerem~$1$, czy
czarną kulę z~numerem~$1$? Mówiąc inaczej, czy $e_{ j } = e_{ 1 }$ lub
$e_{ j } = e_{ 6 }$? To pytanie kieruje nas w~stronę teoriomnogościowego
sformułowania rachunku prawdopodobieństwa. Problem ten możemy bowiem w~prosty
sposób sformułować w~jej języku. Oznaczmy bowiem
$A = \{ e_{ 1 }, e_{ 2 } \}$. Wówczas nasz problem redykuje~się do pytanie
czy $e_{ j } \in A$? Język teorii mnogości okazuje~się wyjątkowo wygodny
w~formułowaniu tego typu pytań, dlatego będziemy go od teraz normalnie
używać.

Problem jaki właśnie rozważyliśmy, można uogólnić w~następujący sposób.
W~tego typu problemie, pytamy nie o~to czy w~danym doświadczeniu
zaszło zdarzenie elementarne takie jak $e_{ 1 }$, tylko czy
to zdarzenie, które zaobserwowaliśmy jest jednym ze~zdarzeń elementarnych
z~pewnego zbioru~$A$. Skoro mówimy o~zbiorach zdarzeń elementarnych, to
teoria mnogości wskazuje, że~powinniśmy sformułować nas problem
w~następujący sposób.

Mamy powtarzalne doświadczenie, w~wyniku wykonania którego daje nam pewien
określony rezultat, każdy taki rezultat nazywamy zdarzeniem elementarnym.
Założyliśmy,
że~w~rozpatrywanym eksperymencie jest tylko skończona liczba możliwych
wyników, więc mamy $n$~zdarzeń elementarnych: $e_{ i }$, $i = 1, 2, 3,\ldots, n$.
Zbiór wszystkich zdarzeń elementarnych nazywamy, zgodnie z~omawianą książką,
\textit{przestrzenią zdarzeń elementarnych} i~oznaczamy~$E$. Teraz
wprowadzamy pojęcia \textit{zdarzenia} (\textit{zdarzenia losowego}),
w~następujący sposób. Niech $A \subseteq E$. Mówimy, że~w~danym eksperymencie
zaszło zdarzenie~$A$, jeśli zaszło zdarzenie elementarne $e_{ j }$
i~$e_{ j } \in A$. Każde zdarzenie elementarne $e_{ j }$ możemy teraz utożsamić
ze zdarzeniem losowym biorąc $A = \{ e_{ j } \}$, pojęcie zdarzenia losowego
jest więc ogólniejsze.

Widzimy teraz, jak intuicja stoi za~pojęciem ogólnego zdarzenia
losowego~$A$. Jest nią pytanie, czy w danym doświadczeniu zaszło jedno ze
zdarzeń elementarnych, które należą do~zbioru~$A$. To jednocześnie tłumaczy
nam na poziomie koncepcyjnym, różnicę między zdarzeniem elementarnym
a~ogólnym zdarzeniem losowym.

Potrzebujemy rozpatrzyć jeszcze co najmniej jeden problem. Wróćmy do naszego
problemu z~urną. Wyciągamy z~niej jedną kulę, następnie ją odkładamy
z~powrotem do urny, z~której ponownie wyciągamy jakąś kulę. Pytamy~się, czy
za pierwszym razem wyciągnęliśmy białą kulę z~numerem $1$, a~za drugim razem
czarną kulę z~numerem~$1$. Inaczej mówiąc, mamy dwuelementowy ciąg wyników
kolejnych wykonań naszego doświadczenia, $a_{ i }$, $i \in \{ 1, 2 \}$ i~zadajemy
sobie pytanie, czy prawdą jest, że~$a_{ 1 } = e_{ 1 }$ i~$a_{ 2 } = e_{ 6 }$?

W~tak postawionym problemie kluczowe jest to, że~uwzględniamy kolejność
wykonania doświadczeń. Jeśli najpierw zaszło zdarzenie elementarne $e_{ 1 }$,
a~potem $e_{ 6 }$ to odpowiedź na wcześniej postawiony problem jest
twierdząca. Jeśli najpierw zaszło zdarzenie $e_{ 6 }$, a~potem $e_{ 1 }$,
czyli najpierw wyciągnęliśmy czarną kulę o~numerze~$1$, a~potem białą
o~numerze $1$, to odpowiedź jest przecząca. Mówiąc inaczej, w~rozważanym
teraz problemie, pytamy~się nie po prostu o~to, czy w~dwóch powtórzenia
doświadczenia zaszły zdarzenie elementarne $e_{ 1 }$ i~$e_{ 6 }$, ale
czy zaszły one i~to w~zadanej kolejności. Problem ten natychmiast~się
uogólnia. Rozpatrzmy dwa zdarzenia
$A_{ 1 }, A_{ 2 } \subseteq E$ i~niech $a_{ i }$, $i = \{ 1, 2 \}$ będzie ciągiem wyników
podwójnego wykonania naszego eksperymentu. Czy jest prawdą,
że~$a_{ i } \in A_{ i }$?

W~ten sposób omówiliśmy najbardziej podstawowe typy problemów jakie
napotykamy w~rachunku prawdopodobieństwa. Po ich uogólnieniu, można je
sprowadzić do dwóch kategorii.



\begin{enumerate}

\item Czy podczas wykonywania doświadczenia zaszło zdarzenie $A$?

\item Czy otrzymany podczas $k$-krotnego powtarzania doświadczenia,
  otrzymaliśmy ciąg zdarzeń losowych $a_{ i }$, o~własności
  $a_{ i } \in A_{ i }$? Gdzie $A_{ i }$ to pewien ciąg zdarzeń losowych,
  $A_{ i } \subseteq E$.

\end{enumerate}



Analizując problem z~rachunku prawdopodobieństwa, należy umieć rozróżnić te
dwie klasy problemów, gdyż pozwala nam to lepiej zrozumieć konkretny
problem, co sprawia, że~jego rozwiązanie staje~się prostsze. Jednocześnie,
z~punktu widzenia formalnego, drugą kategorię możemy sprowadzić do pierwszej,
za pomocą następującego postępowania.

Niech $F^{ k }( E )$ oznacza zbiór wszystkich ciągów o~długości
$k$~o~wartościach w~$E$. Zbiór ten traktujemy jako przestrzeń zdarzeń
elementarny, problemu dla którego pojedyncze doświadczenie polega na
$k$-krotnym wykonaniu jakiegoś powtarzalnego eksperymentu. Należy tutaj
zauważyć, że~pojęcie „powtarzalnego doświadczenia” nie wymusza na nas tego,
że~takie doświadczenie ma~się składać z~pojedynczego działania, co oznacza,
że~również zdarzenie elementarne nie musi być wynikiem pojedynczej czynności.

By to uczynić bardziej zrozumiały, wróćmy do przykładu z~urną. Jeśli
wyciągamy kulę z~urny, wkładamy ją z~powrotem
i~ponownie wyciągamy jakąś kulę, to w~zależności od tego jak chcemy opisać
tą sytuację, możemy mówić, że dwukrotnie wykonano jedno doświadczenie lub
że~jednokrotnie wykonano doświadczenie polegające na dwukrotnym losowaniu
kuli z~urny. Inaczej mówiąc, w~zależności od tego jak opisujemy to co~się
stało, to mamy do czynienia z~dwuelementowym ciągiem zdarzeń elementarnych
lub \textit{jednym} zdarzeniem elementarnym, które polega na dwukrotnym
wyciąganiu kuli z~urny.

Ponieważ te kwestie mogą być źródłem wielu nieporozumień, przyjrzymy~się
im dokładniej. Problemy w~rachunku prawdopodobieństwa mają najczęściej taką
postać. Mamy jakieś powtarzalne doświadczenie, takie jak wyciąganie kuli
z~urny. Już teraz widzieliśmy, że~za pomocą pojedynczej urny z~kulami, możemy
postawić wiele problemów, w~zależności od tego, czy interesuje nas
wyciągnięcie jednej kuli, dwóch kul w~odpowiedniej kolejność,~etc. Z~tego
powodu pojęcie „powtarzalnego doświadczenia” musi być na tyle pojemne,
żeby mogło objąć wszystkie te przypadki. Dlatego możemy zdefiniować takie
powtarzalne doświadczenie, jako wykonanie jednej czynności, jak wyciągnięcie
kuli z~urny, lub wykonanie dwóch, czy więcej czynności, jak wyciągnięcie
dwóch kul z~urny.

Mówiąc inaczej, rachunek prawdopodobieństwa daje nam dużo swobody w~ustaleniu
co jest, a~co nie jest zdarzeniem elementarnym w~danym problemie. To należy
uważać, za zaletę tego podejścia. Problemem leży gdzie indziej. Mianowicie
mając jakąś sytuację, którą chcemy zanalizować za pomocą metod rachunku
prawdopodobieństwa, musimy umieć poprawnie sformułować problem, który
nas interesuje i~następnie zidentyfikować jego zdarzenia elementarne.
Tutaj brak ogólnych reguł postępowania i~każdy nowy przypadek wymaga nowej
analizy i~przebadania. Książki do rachunku prawdopodobieństwa pełne są
różnego typu doświadczeń (problemów), które zostały już zbadane
i~prezentują, lepiej lub gorzej, jak je rozwiązać.

Na stronie $63$ tej książki Hellwig stara~się wyjaśnić, czemu w~książkach
do rachunku prawdopodobieństwa, tak wiele miejsca poświęcono grom losowym.
Fakt, że~sformułowanie konkretnych problemów dotyczących takich gier,
a~następnie zidentyfikowanie zdarzeń elementarnych tychże problemów, jest
zagadnieniem gruntownie przebadanym, jest na pewno kolejnym argument na
rzecz tego, iż~wykłady rachunku prawdopodobieństwa powinny je omawiać.

Problemy należące do obu wymienionych wyżej kategorii można uogólnić na
różne sposoby, my musimy wspomnieć przynajmniej o~jednym z~nich. Niech będą
dane dwa doświadczenia, oznaczmy je $D_{ 1 }$ i~$D_{ 2 }$, zaś ich
przestrzenie zdarzeń elementarnych przez $E_{ 1 }$ i~$E_{ 2 }$. Możemy teraz postawić pytanie, czy wykonując najpierw doświadczenie, a~potem drugie,
dostaniemy ciąg wyników $a_{ i }$, $i \in \{ 1, 2 \}$, taki
że~$a_{ i } \in A_{ i }$, $A_{ i } \subseteq E_{ i }$?

By pokazać, jak naturalny jest tego typu problem, rozważmy następujące
problem. Z~omawianej uprzednio urny wyciągamy kulę i~odkładamy ją na bok,
następnie wyciągamy drugą kulę. Postawy teraz pytanie, czy za pierwszym
razem wyciągnęliśmy białą kulę, a~za drugim czarną? Ponieważ
gdy wyciągamy pierwszą kulę w~urnie jest pięć białych kul i~siedem czarnych,
podczas gdy podczas wyciągania drugiej kuli w~urnie jest $11$ kul, przy czym
liczba kul białych i~czarnych zależy od tego jaką kulę wyciągnęliśmy
poprzednio, to widzimy, że~wyciągnięcie drugiej kuli nie jest prostym
powtórzeniem pierwszej czynności i~możemy całą sytuację opisać, jako dwa
następujące po sobie doświadczenia.

Rozważając ten problem, zadaliśmy pytanie, czy za pierwszym razem
wyciągnęliśmy kulę białą, a~za drugim czarną, więc interesuje nas też
kolejność wyciągania tych kul. Gdybyśmy szukali odpowiedzi na pytanie „Czy
obie kule są czarne?” kolejność wyciągania nie miałaby znaczenia.
Zauważmy, że w~przypadku, gdy kolejność uzyskania wyników musi zostać wzięta
pod uwag, możemy przestrzeń zdarzeń całego elementarnych problemu opisać
jako iloczyn kartezjański przestrzeni zdarzeń elementarnych pierwszego
i~drugiego doświadczenia: $E = E_{ 1 } \times E_{ 2 }$. Dlatego w~takiej sytuacji
możemy mówić o~\textbf{iloczynie kartezjańskim doświadczeń}
$D_{ 1 } \times D_{ 2 }$. Temat iloczynu kartezjańskiego doświadczeń wymaga
głębszej dyskusji, której jednak na razie nie będziemy podejmować.

Jednak taka procedura nie przyda nam~się specjalnie w~analizie problemu
wyciągania dwóch kul, gdyż od tego jaką kulę wyciągnęliśmy za pierwszym
razem, zależy jakie kule możemy wyciągnąć za drugim razem. Mówiąc inaczej,
przestrzeń zdarzeń elementarnych drugiego doświadczenia, zależy od wyniku
pierwszego. Wszystko te rozważania ilustrując, jak trudne potrafi być
pełne sformalizowanie nawet bardzo prostych do opisania słowami
eksperymentów. Do problemów tych może wrócimy w~przyszłości.

Jak można zauważyć, do tej pory nie wspomnieliśmy nic o~prawdopodobieństwie
zdarzeń, co osobom zaznajomiony z~rachunkiem prawdopodobieństwa może
wydawać~się osobliwą decyzją. Uczyniliśmy to jednak z~premedytacją,
chcieliśmy bowiem skupić~się na wyjaśnieniu sensu takich pojęć jak zdarzenie
elementarne, czy przestrzeń zdarzeń elementarnych i~wyjaśnić jak problem
przeprowadzania powtarzalnego doświadczenia motywuje ich wprowadzenie. Jak
widać, do dyskusji tego zagadnienia wystarczy zdolność powiedzenia, czy
doświadczenia dało dany rezultat czy nie, pojęcie prawdopodobieństwa nie
jest do tego w~ogóle potrzebne. Jednocześnie taka analiza pokazuje, dlaczego
język teorii mnogości jest tak wygodny w~rachunku prawdopodobieństwa.

Musimy też zauważyć kolejną ważną rzecz. Z~formalnego punktu widzenia możemy
teraz odrzucić myślenie o~przeprowadzaniu doświadczeń, zostawiają tylko
matematyczny rdzeń, czyli ograniczyć nasze rozważania do operowania
na określonych zbiorach, ich podzbiorach oraz ciągach elementów należących
do tych zbiorów. Jakkolwiek jest
to możliwe, jest to wysoce niewskazane, bowiem gubimy zupełnie z~oczu, czemu
w~ogóle interesują nas takie właśnie problemy. Co gorsza, gdy już
wprowadzimy pojęcie prawdopodobieństwa, to bez pamiętania o~problemie
przeprowadzania powtarzalnego eksperymentu, nie będziemy mogli zrozumieć,
czemu prawdopodobieństwo ma posiadać takie, a~nie inne własności. To zaś
może skutkować tylko zamętem i~utrudnieniem nauki tego działu matematyki.

W~tym momencie potrzebujemy zrobić kolejny krok i~wprowadzić pojęcie
prawdopodobieństwa do rozważanego formalizmu. Choć w~książce Hellwiga
pojęcie to jest dyskutowane trochę dalej i~to w~interesujący sposób, to my
przedyskutujemy je już teraz. Inaczej bowiem nasz rozważania byłby bardzo
niepełne.

Powróćmy do problemu wyciągania kul z~urny. Ponieważ kul czarnych jest
siedem, a~kul białych tylko pięć, intuicyjnie czujemy, iż~wyciągnięcie kuli
czarnej powinno być bardziej prawdopodobne, niż wyciągnięcie kuli białej.
Pojęcie prawdopodobieństwa ma właśnie służyć matematycznemu doprecyzowaniu,
co to znaczy „bardziej prawdopodobne”. Nie będziemy przy tym dyskutować
problemu, czy zawsze takie doprecyzowanie jest możliwe, gdyż to zagadnienie
można omawiać, dopiero po zapoznaniu~się z~podstawami rachunku
prawdopodobieństwa.

Dla problemu z~urną, prawdopodobieństwo jest funkcją
$P : \text{Set}( E ) \to \Rbb_{ + }$, gdzie $\text{Set}( E )$ to zbiór
potęgowy przestrzeni zdarzeń elementarnych,
a~$\Rbb = \{ x \in \Rbb \, | \, x \geq 0 \}$. Na tę funkcję narzucamy szereg
warunków, takich jak to, że~dla każdego
$A \in \text{S}( E )$, $0 \leq P( A ) \leq 1$, jednak w~tym momencie te kwestie nie
są dla nas kluczowe. Jeśli $P( A_{ 1 } ) < P( A_{ 2 } )$ to mówimy,
że~zdarzenie $A_{ 2 }$ jest bardziej prawdopodobne, niż zdarzenie $A_{ 1 }$.
Intuicyjnie odpowiada to temu,
że~jeśli będziemy powtarzać nasze doświadczenie, to zdarzenia elementarne
należące do zdarzenia $A_{ 2 }$ będziemy widzieć częściej, niż te należące
do $A_{ 1 }$.

W~przypadku naszej urny, jeśli będziemy ciągle wyjmować z~niej
kulę i~ponownie ją wkładać, to oczekujemy, iż~częściej wyciągniemy kulę
czarną niż białą. Osoby znające już trochę rachunek prawdopodobieństwa,
wiedzą już, że~standardowa metod rozwiązywania tego zadania pozwala
stwierdzić, że~prawdopodobieństwo wyciągnięcia czarnej kuli to
$\frac{ 7 }{ 12 }$, zaś kuli białej $\frac{ 5 }{ 12 }$.

Bez większego ryzyka można stwierdzić, że~w~rachunku prawdopodobieństwa,
dużo bardziej interesuje nas funkcja prawdopodobieństwa i~jej własności,
niż zdarzenia elementarne i~to tej funkcji poświęcamy najwięcej uwagi.
Przyczyna tego jest dość prosta. Po~wykonaniu doświadczenia możemy
powiedzieć, czy dane zdarzenie zaszło czy nie, za to prawdopodobieństwo
pozwala nam coś powiedzieć, jak bardzo należy~się spodziewać zajścia tego
wydarzenia w~przyszłości, co oczywiście jest dla ludzi bardzo interesujące.
Jednak w~prezentowanym tu podejściu do prawdopodobieństwa, dopiero po
ustaleniu przestrzeni zdarzeń elementarnych jest sens mówienia o~funkcji
prawdopodobieństwa, dlatego tak szczegółowo dyskutowaliśmy te zagadnienia.

Rozwiązując dany problem za pomocą metod rachunku prawdopodobieństwa,
musimy najpierw znaleźć przestrzeń zdarzeń elementarnych, co jest równoważne
określeniu wszystkich zdarzeń elementarnych. Następnie potrzebujemy
zdefiniować wszystkie możliwe zdarzenia. Podczas lektury tej książki
zobaczymy, że~w~przypadku doświadczenia, które ma tylko skończoną liczbę
możliwych wyników, czyli przestrzeń zdarzeń elementarnych~$E$ ma skończoną
moc, każdy
podzbiór~$E$ można rozumieć jako zdarzenie. W~przypadku przestrzeni zdarzeń
elementarnych o~nieskończonej mocy problem ten staje~się znacznie bardziej
skomplikowany i~jest zbyt obszerny by~móc go przedyskutować w~tym miejscu.
Było to jedną z~przyczyn dla których ograniczyliśmy naszą uwagę, do
eksperymentu, który ma tylko skończoną liczbę rezultatów.

Nawet gdy przestrzeń $E$ i~wszystkie zdarzenia zostały już zidentyfikowane,
pozostaje przed nami najtrudniejszy problem: określenie wartości jakie
funkcji prawdopodobieństwa~$P$ przyjmuje na poszczególnych zdarzeniach.
Zadanie to już nawet w~przypadku przestrzeni próbek o~skończonej mocy nie
jest wcale elementarne i~w~dalszym ciągu książki będziemy~się z~nim zmagać
wielokrotnie. Na tym stwierdzeniu możemy w~zasadzie zakończyć naszą analizę
problemu.

W~tych rozważaniach często odnosiliśmy~się do intuicji, więc pewne pojęcia
są w~naturalny sposób niedookreślone. Wydaje~się nam, że~na tę chwilę
podstawowe pojęcia rachunku prawdopodobieństwa powinny być na tyle
zrozumiałe, żeby można było na ich podstawie przejść dalej.

\VerSpaceFour





\noindent
\Str{36} W~tym miejscu autor omawia problem dwukrotnego rzutu monetą,
by zilustrować, że~iloczyn zdarzeń nie musi oznaczać, iż~oba zdarzenia
zaszły jednocześnie. Zamieszczanie tego typu przykładów, ma niewątpliwie na
celu nauczenie czytelnika, jak konkretne doświadczenia powtarzalna można
opisać za pomocą pojęć rachunku prawdopodobieństwa, to zaś jest wstępem
do rozwiązania danego problemu za pomocą odpowiednich metod tegoż działu
matematyki. W~tym konkretnym wypadku, przykład ma służyć wyjaśnieniu,
że~interpretacja wedle której, iloczyn zdarzeń oznacza, iż~dwa zdarzenia
musiały zajść w~tym samym czasie, jest błędna.

Wobec tego, zajmujemy~się tutaj dość subtelnym zagadnieniem: jak należy
rozumieć i~stosować poznane pojęcia rachunku prawdopodobieństwa do opisu
konkretnego problemu? By uczynić to łatwiejszym do zrozumienie,
przeanalizujemy teraz problem podwójnego rzutu monetą.

Pojedynczy rzut monetą jest doświadczenie powtarzalnym. Jak to było
już wcześniej omówione, biorąc za punkt wyjścia takie doświadczenie możemy
zbadać wiele różnych problemów. W~tym wypadku możemy oczywiście analizować
pojedynczy rzutu, dwa rzuty następujące po sobie, trzy rzutu następujące po
sobie,~etc. My ograniczymy~się do problemu dwóch rzutów monetą. Autor, choć
nie napisał tego jawnie, przyjął, że w~tym kontekście przez
\textit{pojedyncze} doświadczenie rozumie on podwójnym rzucie monetą. Dla
uniknięcia niejasności, przez podwójny rzut monetą rozumiemy sytuację,
w~której najpierw rzucamy monetą, a~potem tę samą monetą rzucamy drugi raz.
Tym samym wykluczamy z~rozważań problem, który powstaje, gdy dwoma monetami
rzucamy w~tej samej chwili. Wyjaśnimy dlaczego tak należy rozumieć ten
fragment książki, po tym jak opiszemy dokładniej problem podwójnego rzutu
monetą.

Wypadnięcie orła przy rzucie monetą będziemy oznaczać przez $\Orzel$,
zaś~wypadnięcie reszki przez $\Reszka$. To, że~przy podwójnym rzucie
najpierw wypadł orzeł, a~potem reszka, będziemy oznaczali przez $\Orzel
\Reszka$. Analogicznie $\Reszka \Orzel$ oznacza, że~najpierw wypadła reszka,
a~potem orzeł. Przestrzeń zdarzeń elementarnych takiego doświadczenia
zawiera~$2^{ 2 } = 4$ elementów, gdyż tyle jest ciągów o~długości dwa,
przyjmujących wartości w~zbiorze dwuelementowym. Możemy więc wypisać
wszystkie elementy tej przestrzeni w~sposób jawny:
\begin{equation}
  \label{eq:Hellwig-Elementy-ETC-01}
  E = \{ \Orzel \Orzel, \Orzel \Reszka, \Reszka \Orzel, \Reszka \Reszka \}.
\end{equation}

Zapytajmy~się teraz, czy można za~pomocą tego co zostało powiedziane opisać
zdarzenie wypadnięcia orła w~pierwszym rzucie? Odpowiedź jest twierdząca,
zdarzenie to opisuje zbiór
\begin{equation}
  \label{eq:Hellwig-Elementy-02}
  E_{ 1 } = \{ \Orzel \Orzel, \Orzel \Reszka \}.
\end{equation}
Od razu widać, że~$E_{ 1 } \subseteq E$. Wypadnięcie orła w~drugim rzucie opisuje
zdarzenie
\begin{equation}
  \label{eq:Hellwig-Elementy-03}
  E_{ 1 } = \{ \Orzel \Orzel, \Reszka \Orzel \}.
\end{equation}
Widzimy, że~teoriomnogościowy iloczyn tych zdarzeń to
\begin{equation}
  \label{eq:Hellwig-Elementy-03}
  E_{ 1 } \cap E_{ 2 } = \{ \Orzel \Orzel \},
\end{equation}
czyli dokładnie to zdarzenie o~które nam chodziło: wypadnięcie orła
w~pierwszym i~drugim rzucie.

Możemy teraz wyjaśnić, czemu ten problem należy rozumieć w~taki sposób,
że~zdarzeniem elementarnym jest ciąg dwóch rzutów monetą. Jeśli bowiem chcemy
opisywać iloczyn zdarzeń za~pomocą teorii mnogości\footnote{Nie będziemy~się
  wdawać w~dyskusję, czy należy sformalizować rachunek prawdopodobieństwa,
  bez użycia języka teorii mnogości.}, to należy każde zdarzenie opisać jako
podzbiór
przestrzeni zdarzeń elementarnych. W~takiej sytuacji zdarzenie,
że~w~pierwszym rzucie wypadł orzeł należy rozumieć jako zbiór wszystkich
możliwych wyników podwójnego rzutu monetą, takich że~w~pierwszym z~nich
wypadł orzeł. Dokładnie to opisuje zbiór $E_{ 1 }$.

Przyjrzyjmy~się teraz zdaniu, jakie możemy znaleźć w~ostatniej linii
omawianej strony: „W~przykładzie tym zdarzenia $E_{ 1 }$ i~$E_{ 2 }$ nie
zachodzą równocześnie.”. Autorowi chodziło o~to, że~gdy rzucam dwa razy
monetą i~w~obu wypadkach wypada orzeł, to pierwsze wypadnięcie orła
rozgrywa~się
wcześniej niż drugie, czyli ewidentnie nie występują one jednocześnie.
Jednak tutaj pojawia~się następująca subtelność. W~prezentowanym tu
formalizmie zdarzenia $E_{ 1 }$ i~$E_{ 2 }$ składają~się z~ciągów wyników
podwójnego rzutu monetą, ciężko więc ich zajściu przypisać \textit{jedną}
wartość czasu.

W~obecnej sytuacji najbardziej rozsądne jest przyjęcie, że~każdemu zdarzeniu
$A \subseteq E$ przyporządkowujemy cały okres czasu, jaki potrzebny jest na
wykonania dwóch rzutów. Wtedy jednak można powiedzieć,
że~$E_{ 1 }$ i~$E_{ 2 }$ występując
jednocześnie, gdy ich okresy trwania są identyczne. Pokazuje to, że~ustalenie
odpowiedniości między intuicyjny pojęciem zdarzenia, a~opartym na teorii
mnogości sformułowaniem rachunku prawdopodobieństwa kryje wiele subtelności
i~pułapek.

W~tym momencie najlepiej jest ten zamęt pojęciowy wyjaśnić w~następujący
sposób. Na omawianej stronie iloczyn zdarzeń jest najpierw ilustrowany przez
problem wyciągania karty z~talli, dopiero potem przez podwójny rzut monetą,
dlatego warto zwrócić uwagę co oba te przykłady mają wspólnego, a~co je
różni. Zacznijmy od problemu wyciągania karty z~tali. Zdarzenie będące
iloczyny wyciągnięcia koloru kier i~zdarzenia polegającego na
wyciągnięciu damy, polega na tym, że~karta którą wyciągamy
jest damą kier. Mówiąc inaczej iloczyn tych dwóch zdarzeń polega na
wyciągnięciu karty, która jest jednocześnie kartą koloru kier i~damą, więc
oba te zdarzenia \textit{muszą} zajść jednocześnie.

W~przypadku podwójnego rzutu monetą myślimy o~eksperymencie, który nie
rozgrywa~się w~jednej chwili, ale potrzebuje dłuższego czasu. W~naszym
wypadku, niech pierwszy rzut monetą odbywa~się w~chwili $t_{ 1 }$, a~drugi
w~chwili~$t_{ 2 }$. Zdarzenie $E_{ 1 }$, wyrzucenie orła w~pierwszym rzucie,
jest wyznaczone przez to co stało~się w~chwili $t_{ 1 }$, zaś zdarzenie
$E_{ 2 }$ przez wyrzucenie orła w~chwili $t_{ 2 }$. Możemy teraz zapytać~się
o~iloczyn zdarzeń $E_{ 1 }$ i~$E_{ 2 }$, pomimo tego, że~zdarzenia
$E_{ 1 }$ i~$E_{ 2 }$ są wyznaczone, przez rzuty monetą, które nie
rozegrały~się w~tej samej chwili.

To jednak pokazuje jaki jest problem ze zdaniem „W~przykładzie tym zdarzenia
$E_{ 1 }$ i~$E_{ 2 }$ nie zachodzą równocześnie.”, gdyż mówi ono, że~zdarzenie
wyrzucenia orła w~pierwszym rzucie i~w~drugim rzucie nie rozgrywają~się
jednocześnie. Nasza analiza powyżej mówi, że~zdarzeniom tym nie można
przypisać jednej chwili czasu, a~jeśli przypiszemy im przedział czasu, to
możemy powiedzieć, że~zachodzą jednocześnie.

Co więcej, problem ten pozostaje aktualny, nawet, gdy nie będziemy używać
formalizmu rachunku prawdopodobieństwa opartego na teorii mnogości.
Jeśli~się bowiem zastanowimy, zdarzenie wyrzucenia orła w~drugim rzucie
oznacza, że~przedtem już wykonaliśmy jakiś rzut monetą i~dał on jakiś wynik.
To samo jest prawdą o~wyrzuceniu orła w~pierwszym rzucie, nawet jeśli
zrozumienie tego jest trochę trudniejsze, niż w~poprzednim przypadku.

\VerSpaceFour





\noindent
\Str{47}






% ######################################
% \CenterBoldFont{Błędy}


% \begin{center}

%   \begin{tabular}{|c|c|c|c|c|}
%     \hline
%     Strona & \multicolumn{2}{c|}{Wiersz} & Jest & Powinno być \\
%     \cline{2-3}
%     & Od góry & Od dołu & & \\
%     \hline
    % & & & & \\
    % & & & & \\
    % & & & & \\
%     \hline
%   \end{tabular}

% \end{center}

% ######################################

% \VerSpaceSix



% ####################################################################







% ####################################################################
% ####################################################################
% Bibliography

\printbibliography





% ############################
% End of the document

\end{document}
